\documentclass[a4,11pt]{aleph-notas}

% -- Paquetes adicionales
\usepackage{enumitem}
\usepackage{aleph-comandos}

% -- Datos
\institucion{Facultad de Ciencias Exactas, Naturales y Ambientales}
\carrera{Catálogo STEM}
\asignatura{Matemática III}
\tema{Resumen 04: Límites y diferenciabilidad}
\autor[A. Merino]{Andrés Merino}
\fecha{Periodo 2026-1}

\logouno[0.16\textwidth]{Logos/logoPUCE_04_ac}
\definecolor{colortext}{HTML}{1957d1}
\definecolor{colordef}{HTML}{1957d1}
\fuente{montserrat}


\begin{document}

\encabezado

En este resumen consideraremos $\Omega\subset \R^n$ e $I\subset\R$ un intervalo.

%%%%%%%%%%%%%%%%%%%%%%%%%%%%%%%%%%%%%%
\section{Límites y continuidad de campos escalares}
%%%%%%%%%%%%%%%%%%%%%%%%%%%%%%%%%%%%%%

\begin{defi}[Límite]
    Dados un campo escalar $\func{f}{\Omega}{\R}$, $a\in \Omega'$ y $L\in\R$, se dice que $L$ es el límite de $f$ cuando $x$ tiende a $a$ si: para todo $\epsilon>0$, existe $\delta>0$ tal que para todo $x\in \Omega$,
    \[
        0<\|x-a\|<\delta \Imp |f(x)-L| < \epsilon.
    \]
    Esto se denota por $\dlim_{x\to a}f(x)=L$.
\end{defi}

\begin{prop}
    El límite de una función es único.
\end{prop}

\begin{defi}[Límite por trayectorias]
    Sea $\func{f}{\Omega}{\R}$ un campo escalar, $a\in \Omega'$, $\func{\alpha}{I}{\R^n}$ una trayectoria inyectiva tal que existe $t_0\in I$ con $\alpha(t_0)=a$. Se define el límite de $f$ cuando $x$ tiende a $a$ por la trayectoria $\alpha$ mediante
    \[
        \lim_{t\to t_0} f(\alpha(t)),
    \]
    siempre que este exista.
\end{defi}

\begin{prop}
    Sean $\func{f}{\Omega}{\R}$ un campo escalar y $a\in \Omega'$. Si existe $\dlim_{x\to a}f(x)=L$, entonces existe el límite por cualquier trayectoria y es igual a $L$.
\end{prop}

\begin{prop}[Propiedades de los límites]
    Sean $\func{f,g}{\Omega}{\R}$ campos escalares y $a\in \Omega'$. Si existen $\dlim_{x\to a}f(x)$ y $\dlim_{x\to a}g(x)$, entonces
    \begin{itemize}
        \item $\dlim_{x\to a}\l(f(x)+g(x)\r)=\dlim_{x\to a}f(x)+\dlim_{x\to a}g(x)$;
        \item $\dlim_{x\to a}(f(x)g(x))=\l(\dlim_{x\to a}f(x)\r)\l(\dlim_{x\to a}g(x)\r)$;
        \item $\dlim_{x\to a}\l(\frac{f(x)}{g(x)}\r)=\frac{\dlim_{x\to a}f(x)}{\dlim_{x\to a}g(x)}$, siempre que $\dlim_{x\to a}g(x)\neq 0$.
    \end{itemize}
\end{prop}

\begin{defi}[Continuidad]
    Dados $\func{f}{\Omega}{\R}$ un campo escalar y $a\in \Omega$, se dice que $f$ es continua en $a$ si
    \[
        \lim_{x\to a}f(x)=f(a).
    \]
    Se dice que $f$ es continua en $\Omega$ si $f$ es continua en cada punto de $\Omega$.
\end{defi}

\begin{prop}
    Dados $\func{f}{\Omega}{\R}$ un campo escalar continuo con $\Omega$ cerrado y acotado, se tiene que $f$ alcanza su mínimo y su máximo en $\Omega$.
\end{prop}


\end{document}
